% !TeX program = xelatex
% !TeX encoding = UTF-8
%% Copyright 2012-2013 by abnTeX2 group at http://abntex2.googlecode.com/ 
% Adaptado para uso no DACC - IF Sudeste MG, campus Rio Pomba
% ABNT NBR 14724:2011: Informacao e documentacao - Trabalhos academicos
% ------------------------------------------------------------------------
% PACOTES
% ------------------------------------------------------------------------
\documentclass[
	% -- opções da classe memoir --
	12pt,				% tamanho da fonte
	openright,			% capítulos começam em pág ímpar (insere página vazia caso preciso)
	oneside,			% para impressão em verso e anverso. Oposto a oneside
	a4paper,			% tamanho do papel. 
	% -- opções da classe abntex2 --
	%chapter=TITLE,		% títulos de capítulos convertidos em letras maiúsculas
	%section=TITLE,		% títulos de seções convertidos em letras maiúsculas
	%subsection=TITLE,	% títulos de subseções convertidos em letras maiúsculas
	%subsubsection=TITLE,% títulos de subsubseções convertidos em letras maiúsculas
	% -- opções do pacote babel --
	english,			% idioma adicional para hifenização
	french,				% idioma adicional para hifenização
	spanish,			% idioma adicional para hifenização
	brazil,				% o último idioma é o principal do documento
	]{abntex2}

\usepackage{cmap}				% Mapear caracteres especiais no PDF
%\usepackage[T1]{fontenc}		% Selecao de codigos de fonte.
%\usepackage[utf8]{inputenc}		% Codificacao do documento (conversão automática dos acentos)
\usepackage{lastpage}			% Usado pela Ficha catalográfica
\usepackage{indentfirst}		% Indenta o primeiro parágrafo de cada seção.
\usepackage{color}				% Controle das cores
\usepackage{graphicx}			% Inclusão de gráficos
\usepackage{underscore}
\usepackage{float}
\usepackage[final]{pdfpages}
\usepackage{multicol}
\usepackage{fontspec}
\usepackage{array}
\usepackage{xltabular}   % longtable + tabularx
\usepackage{booktabs}
%\setmainfont{Arial}
% Adicionado Manualmente
\usepackage{fontspec}
\setmainfont{Liberation Sans}
\usepackage{lipsum}				% para geração de dummy text
\usepackage[brazilian,hyperpageref]{backref}	 % Paginas com as citações na bibl
\usepackage[alf]{abntex2cite}	% Citações padrão ABNT
\renewcommand{\backrefpagesname}{}
\renewcommand{\backref}{}




% ------------------------------------------------------------------------
% Informações de dados para CAPA e FOLHA DE ROSTO
% ------------------------------------------------------------------------

\titulo{Desenvolvimento Full Stack de um simulador de banco para trabalhar matemática financeira com estudantes do ensino médio.}


\autor{Henrique Casagrande Bravin}
\local{Manhuaçu}
\data{2026}
\orientador{Prof. Ms.  ROBSON DE SOUZA RESENDE}
\coorientador{Nome do Coorientador}
\tipotrabalho{Trabalho de Conclusão de Curso}
\preambulo{Trabalho de Conclusão curso  apresentado ao \textit Campus Manhuaçu, do Instituto Federal de Educação, Ciência e Tecnologia do Sudeste de Minas Gerais, como parte das exigências do curso de Pós-Graduação Lato Sensu em Desenvolvimento Full Stack.}
% ---


% ---
% Configurações de aparência do PDF final

% alterando o aspecto da cor azul
\definecolor{blue}{RGB}{41,5,195}

% informações do PDF
\makeatletter
\hypersetup{
     	%pagebackref=true,
		pdftitle={\@title}, 
		pdfauthor={\@author},
    	pdfsubject={\imprimirpreambulo},
	    pdfcreator={Nome do aluno(a)}, %procurar
		pdfkeywords={content-based image retrieval}{desenvolvimento web}{exame de fundo de olho}{histograma backprojection}{íris}, 
		colorlinks=true,       		% false: boxed links; true: colored links
    	linkcolor=black,          	% color of internal links
    	citecolor=black,        		% color of links to bibliography
    	filecolor=black,      		% color of file links
		urlcolor=black,
		bookmarksdepth=4
}
\makeatother
% --- 

% --- 
% Espaçamentos entre linhas e parágrafos 
% --- 

% O tamanho do parágrafo é dado por:
\setlength{\parindent}{1.3cm}

% Controle do espaçamento entre um parágrafo e outro:
\setlength{\parskip}{0.2cm}  % tente também \onelineskip

% ---
% compila o indice
% ---
\makeindex
% ---

% ----
% Início do documento
% ----
\begin{document}

% Retira espaço extra obsoleto entre as frases.
\frenchspacing 

% ----------------------------------------------------------
% ELEMENTOS PRÉ-TEXTUAIS
% ----------------------------------------------------------
% \pretextual

% ---
% Capa
% ---
\begin{center}
\textbf{ 
INSTITUTO FEDERAL DE EDUCAÇÃO, CIÊNCIA E TECNOLOGIA DO SUDESTE DE MINAS GERAIS - CAMPUS RIO POMBA}
\end{center}

\imprimircapa
% ---

% ---
% Folha de rosto
% (o * indica que haverá a ficha bibliográfica)
% ---
\imprimirfolhaderosto*
% ---

% ---
% Inserir a ficha bibliografica
% ---

% Para inserir a ficha catalográfica oficial da biblioteca, utilize:
% \begin{fichacatalografica}
%     \includepdf{fig_ficha_catalografica.pdf}
% \end{fichacatalografica}

% Para inserir a ficha catalográfica temporária, utilize:

\begin{fichacatalografica}
	\vspace*{\fill}					% Posição vertical
	\hrule							% Linha horizontal
	\begin{center}					% Minipage Centralizado
	\begin{minipage}[c]{12.5cm}		% Largura
	
	FICHA CATALOGRÁFICA TEMPORÁRIA \\
	\imprimirautor
	
	\hspace{0.5cm} \imprimirtitulo  / \imprimirautor. --
	\imprimirlocal, \imprimirdata-

	\hspace{0.5cm} \imprimirorientadorRotulo~\imprimirorientador\\
	
	\hspace{0.5cm}
	\parbox[t]{\textwidth}{\imprimirtipotrabalho~--~Instituto Federal de Educação, Ciência e Tecnologia do Sudeste de Minas, Campus Rio Pomba,
	\imprimirdata.}\\

	
	\end{minipage}
	\end{center}
	\hrule
\end{fichacatalografica}
% ---

% ---
% Inserir errata
% ---
%\begin{errata}
%Elemento opcional da \citeonline[4.2.1.2]{NBR14724:2011}. %Exemplo:

%\vspace{\onelineskip}
%
%FERRIGNO, C. R. A. \textbf{Tratamento de neoplasias ósseas apendiculares com
%reimplantação de enxerto ósseo autólogo autoclavado associado ao plasma
%rico em plaquetas}: estudo crítico na cirurgia de preservação de membro em
%cães. 2011. 128 f. Tese (Livre-Docência) - Faculdade de Medicina Veterinária e
%Zootecnia, Universidade de São Paulo, São Paulo, 2011.

%\begin{table}[htb]
%\center
%\footnotesize
%\begin{tabular}{|p{1.4cm}|p{1cm}|p{3cm}|p{3cm}|}
%  \hline
%   \textbf{Folha} & \textbf{Linha}  & \textbf{Onde se lê} % & \textbf{Leia-se}  \\
%    \hline
%    1 & 10 & auto-conclavo & autoconclavo\\
%   \hline
%\end{tabular}
%\end{table}
%
%\end{errata}
% ---

% ---
% Inserir folha de aprovação
% ---
% Se possuir a folha de aprovação emitida pela biblioteca, utilizar o comando abaixo para inserir o PDF. Senão, use o modelo temporário abaixo
% \includepdf{folhadeaprovacao_final.pdf}
%
\begin{folhadeaprovacao}

  \begin{center}
    {\ABNTEXchapterfont\large\imprimirautor}

    \vspace*{\fill}\vspace*{\fill}
    {\ABNTEXchapterfont\bfseries\Large\imprimirtitulo}
    \vspace*{\fill}
    
    \hspace{.45\textwidth}
    \begin{minipage}{.5\textwidth}
        \imprimirpreambulo
    \end{minipage}%
    \vspace*{\fill}
   \end{center}
    
   Trabalho aprovado. \imprimirlocal, 25 de novembro de 2020:

   \assinatura{\textbf{\imprimirorientador} Orientadora, IF Sudeste MG - Rio Pomba} 
   \assinatura{\textbf{Nome do Coorientador}, \\ Coorientador, IF Sudeste MG - Rio Pomba}
   \assinatura{\textbf{PROFESSOR X} \\ IF Sudeste MG - Rio Pomba}
   \assinatura{\textbf{PROFESSOR Y} \\ IF Sudeste MG - Rio Pomba}
      
   \begin{center}
    \vspace*{0.5cm}
    {\large\imprimirlocal}
    \par
    {\large\imprimirdata}
    \vspace*{1cm}
  \end{center}
  
\end{folhadeaprovacao}
% ---

% ---
% Dedicatória
% ---
\input{preTextual/dedicatoria}
% ---

% ---
% Agradecimentos
% ---
\input{preTextual/agradecimentos}
% ---

% ---
% Epígrafe
% ---
\input{preTextual/epigrafe}
% ---

% ---
% RESUMOS
% ---

% resumo em português
\input{preTextual/resumoBR}

% resumo em inglês
\input{preTextual/resumoUS}


% ---
% inserir lista de ilustrações
% ---
\pdfbookmark[0]{\listfigurename}{lof}
\listoffigures*
\cleardoublepage
% ---

% ---
% inserir lista de tabelas
% ---
\pdfbookmark[0]{\listtablename}{lot}
\listoftables*
\cleardoublepage
% ---

% ---
% inserir lista de abreviaturas e siglas
% ---
\input{Textual/listaSiglas}
% ---

% ---
% inserir lista de símbolos
% ---
%\begin{simbolos}
%  \item[$ \Lambda $] Lambda
%\end{simbolos}
% ---

% ---
% inserir o sumario
% ---
\pdfbookmark[0]{\contentsname}{toc}
\addcontentsline{arquivo}{unidade}{entrada}
\tableofcontents*
\cleardoublepage
% ---



% ---------------------------------------------------------------------------------------------
% ELEMENTOS TEXTUAIS
% ---------------------------------------------------------------------------------------------
\textual
\setcounter{page}{1}

\chapter{Introdução}
\label{cap-introducao}
\section{Contextualização e justificativa}
\label{sec-introducao-contextualizacao}

Pagamentos, compras parceladas e operações de crédito exigem a
interpretação de informações como saldo disponível, taxa de juros,
prazo e valor total a pagar. Ao investigar a alfabetização
financeira, Potrich, Vieira e Kirch consideram três dimensões:
conhecimento, atitude e comportamento financeiros
\cite{potrich2015}. No contexto escolar, uma situação de compra
permite articular o cálculo à decisão que ele informa. Saber
determinar o valor de uma prestação, por exemplo, é parte da
análise necessária para decidir se a compra cabe no orçamento.

Neste trabalho, apresenta-se o EDBank, um simulador bancário
desenvolvido para apoiar atividades de matemática financeira no
Ensino Médio. O participante utiliza uma conta fictícia vinculada
a uma agência, na qual pode receber um salário simulado, pagar
cobranças e realizar compras. Como essas operações afetam a mesma
conta, uma escolha feita durante a atividade modifica os recursos
disponíveis para a decisão seguinte.

O trabalho descreve como esse ambiente foi especificado e
implementado. O desenvolvimento compreende a definição das regras
das operações, a organização dos dados e a integração entre
interface e servidor. A finalidade educacional orienta a escolha
das situações representadas pelo simulador.

\section{Objetivo geral}
\label{sec-introducao-objetivo-geral}

Desenvolver uma aplicação web full stack que simule operações bancárias fictícias para apoiar o ensino de matemática financeira no Ensino Médio.


\section{Objetivos Específicos}
\label{sec-introducao-objetivo-especifico}

\begin{itemize}
    \item Especificar os requisitos funcionais e as regras de negócio das operações previstas para o simulador.

    \item Projetar a arquitetura da aplicação e modelar os dados necessários ao registro de contas e movimentações fictícias.

    \item Implementar os fluxos de abertura de conta, consulta de saldo e extrato, pagamentos, transferências e compras parceladas no ambiente simulado.

    \item Verificar o funcionamento das operações implementadas por meio de cenários de teste que contemplem cálculos financeiros, atualização de saldos e tratamento de operações inválidas.
\end{itemize}

\chapter{Fundamentação Teórica}
\label{cap-fundamentacao}
\section{Educação financeira e matemática financeira no Ensino Médio}
\label{sec-fundamentacao-educacao-financeira}

A educação financeira abrange a compreensão de situações que envolvem o uso do dinheiro, o planejamento de gastos e a avaliação de decisões relacionadas ao consumo e ao crédito. A matemática financeira fornece instrumentos para analisar essas situações, como o cálculo de porcentagens e juros e a comparação de valores ao longo do tempo. O Banco Central do Brasil aborda a gestão do orçamento e o uso do crédito como temas da formação financeira \cite{bcb2026caderno}.


No Ensino Médio, essa relação encontra respaldo na Base Nacional
Comum Curricular. A habilidade EM13MAT203 contempla o uso de
conceitos matemáticos em atividades de planejamento e análise
com aplicativos e planilhas, incluindo situações de orçamento
e simulação de juros \cite{brasilBNCC2018}. Para o trabalho em
sala, essa orientação permite formular tarefas em que o estudante
escolhe entre alternativas, registra os cálculos realizados e
justifica sua decisão com base nos valores obtidos.


\section{Operações bancárias e movimentação de saldo}
\label{sec-fundamentacao-operacoes}

O saldo expressa a posição da conta em determinado instante.
O extrato registra as movimentações que explicam a passagem de
um saldo a outro. Se \(S_0\) representa o saldo inicial e \(m_k\)
representa o valor da \(k\)-ésima movimentação, com sinal positivo
para entradas e negativo para saídas, o saldo após \(n\)
movimentações é dado por
\[
    S_n = S_0 + \sum_{k=1}^{n} m_k.
\]

Por exemplo, uma conta com saldo inicial de \mbox{R\$~200,00}
que registra um pagamento de \mbox{R\$~45,00} e recebe uma
transferência de \mbox{R\$~30,00} passa a apresentar saldo de
\mbox{R\$~185,00}. A leitura do extrato permite reconstruir esse
resultado e explicar por que o saldo diminuiu em
\mbox{R\$~15,00}, embora a conta tenha recebido um crédito.

\section{Juros, prazo e comparação de valores}
\label{sec-fundamentacao-juros}



Em uma operação de crédito, o valor da prestação informa quanto
será pago em cada vencimento. Para examinar o custo da decisão,
é preciso considerar também o número de pagamentos, o valor
financiado e os encargos incluídos. Em um financiamento com
\(n\) prestações iguais de valor \(P\), sem entrada ou cobranças
adicionais, o total pago corresponde a \(nP\). Duas propostas
podem, portanto, apresentar prestações de valores próximos e
totais distintos quando seus prazos diferem. A metodologia da
Calculadora do Cidadão relaciona prestação, prazo, taxa mensal
e valor financiado em operações com capitalização mensal
\cite{bcbMetodologiaFixas}.

Nas aplicações financeiras, a questão de interesse é a evolução
do capital investido. Como referência matemática, um capital
inicial \(C\), aplicado a uma taxa composta constante \(i\) por
\(n\) períodos, sem aportes, retiradas ou encargos, resulta no
montante \(M=C(1+i)^n\). A unidade do período deve corresponder
à periodicidade da taxa. Uma taxa mensal, por exemplo, exige que
o prazo seja expresso em meses.

Esses conceitos orientam perguntas distintas em uma atividade
escolar. No crédito, o estudante pode investigar quanto pagará
ao longo do prazo escolhido. No investimento, pode examinar como
o valor aplicado evolui durante a atividade. No EDBank, a leitura
dos resultados de investimento deve considerar a modalidade
simulada e as regras de rendimento adotadas pelo sistema.


\section{Simulação digital como recurso didático}
\label{sec-fundamentacao-simulacao}

Em uma atividade de matemática financeira, a simulação permite
alterar uma condição da operação e examinar seus efeitos sobre o
resultado. Ao comparar dois cenários, o estudante precisa identificar
quais valores foram mantidos, qual parâmetro foi modificado e como
essa mudança explica a diferença observada.

Uma experiência brasileira com essa abordagem é apresentada por
Lima e Batista em um guia didático elaborado a partir de uma
intervenção com 33 estudantes do Ensino Médio de uma escola estadual
de Alagoas. A proposta articula o registro de despesas, a análise de
padrões de consumo e a simulação de aportes periódicos. Em uma das
etapas, os estudantes utilizam os dados produzidos pelo aplicativo
para fundamentar suas respostas. A interpretação dos resultados
integra, assim, a própria atividade de matemática financeira
\cite{limaBatista2026}.

No EDBank, uma atividade com finalidade semelhante pode partir da
compra de um mesmo produto em três ou quatro parcelas. O professor
pode pedir que os estudantes antecipem o total a pagar em cada
condição, consultem os valores apresentados pelo simulador e
expliquem a diferença a partir da regra de acréscimo adotada pelo
shopping. A tarefa relaciona porcentagem, número de parcelas e
custo total em uma situação que o estudante pode reproduzir.

\chapter{Soluções Relacionadas}
\label{cap-solucoes}
\section{Calculadora do Cidadão}
\label{cap-solucoes-calculadora}

A Calculadora do Cidadão é uma ferramenta disponibilizada pelo Banco Central do Brasil para simular situações do cotidiano financeiro. Na modalidade de financiamento com prestações fixas, a ferramenta permite relacionar o valor financiado, a taxa mensal de juros, o número de meses e o valor da prestação. O usuário informa os dados conhecidos para calcular a variável de interesse \cite{bcbCalculadora}.


Para este projeto, a ferramenta serve como referência para a apresentação das informações de uma simulação de crédito. A identificação da taxa mensal, do prazo e do valor financiado ajuda o usuário a reconhecer quais dados foram utilizados para obter o resultado.

\section{Comparação com a proposta do simulador bancário}
\label{cap-solucoes-comparacao}

A Calculadora do Cidadão oferece uma referência para estruturar a interação entre os dados informados e o resultado de um cálculo financeiro \cite{bcbCalculadora}. O simulador proposto neste trabalho situa esse tipo de cálculo em um ambiente didático que também representa contas, saldos e movimentações fictícias.


A aplicação prevê dois tipos de atividade. Na primeira, o estudante realiza pagamentos e transferências e verifica seus efeitos sobre o saldo e o extrato. Na segunda, altera parâmetros de uma simulação de crédito e compara os resultados obtidos. Essas atividades permitem trabalhar tanto a interpretação de movimentações quanto a análise de taxa, prazo e custo.

\chapter{Concepção do Simulador}
\label{cap-simulador}
O EDBank é uma aplicação web que simula operações bancárias para uso
em atividades de matemática financeira no Ensino Médio. A aplicação
organiza as contas em agências e permite aos participantes movimentar
valores fictícios, acompanhar registros no extrato e consultar as
consequências financeiras de suas operações.

A concepção do sistema reúne dois tipos de interação. Na área da
conta, o participante realiza operações como pagamentos,
transferências e compras. Na área de gerência, é possível configurar
a agência, acompanhar suas contas e executar ciclos de processamento.
Entre as operações realizadas em um ciclo está o crédito de salário
fictício, associado à profissão atribuída no momento da abertura da
conta.

\section{Público-alvo e Contexto de Uso}
\label{cap-cencepcao-publico}

O EDBank destina-se a estudantes do Ensino Médio em atividades de
matemática financeira. Em uma atividade orientada pelo professor,
o estudante pode prever o saldo após uma sequência de operações,
conferir o resultado no extrato e explicar como os pagamentos, os
créditos e as transferências afetaram o valor disponível.

A aplicação oferece uma área de conta para os participantes e uma
área de gerência para administrar agências. Também há previsão de
subgerência por convite e uma área administrativa do sistema.
Em um possível cenário de uso em sala de aula, o professor pode
utilizar a gerência de uma agência para organizar a atividade e
acompanhar as contas vinculadas a ela. A agência é a unidade de
organização adotada pela aplicação.

\section{Escopo da aplicação}
\label{cap-concepcao-escopo}

O EDBank compreende a abertura de contas fictícias em agências,
a consulta de saldo e extrato, as transferências entre contas,
o pagamento de cobranças, o uso de cartões, os investimentos
simulados e as compras em um shopping interno. Na área de gerência,
o sistema reúne funções relacionadas à configuração da agência,
à administração das contas e ao processamento de ciclos.

As movimentações utilizam valores fictícios e são processadas no
ambiente da aplicação, sem liquidação financeira real. A
autenticação com conta Google é utilizada para identificar o
participante no acesso ao sistema.

O módulo de empréstimos tem caráter demonstrativo. O cálculo está
previsto na interface, e a ação de solicitação exibe um protocolo
no navegador. Na versão analisada, essa ação não gera contrato,
registro de parcelas ou crédito em conta. A interface também inclui
recursos relacionados ao Pix, mas o fluxo de envio está incompleto
na versão analisada.


\section{Requisitos funcionais}
\label{cap-concepcao-requisitos-funcionais}

Os requisitos funcionais a seguir descrevem as ações previstas
para os participantes e para a gerência, a partir dos fluxos
identificados no código da aplicação.

\renewcommand{\arraystretch}{1.4}

\begin{xltabular}{\textwidth}{|>{\centering\arraybackslash}m{1.6cm}|>{\centering\arraybackslash}m{3.6cm}|X|}
\caption{Requisitos Funcionais do Sistema}\label{tab:requisitos-funcionais}\\
\hline
\textbf{Código} & \textbf{Item} & \textbf{Descrição} \\
\hline

RF01 & Autenticação & Permitir o acesso do participante por meio de uma conta Google. \\
\hline
RF02 & Abertura de conta & Permitir que o participante abra uma conta fictícia em uma agência, mediante código de acesso quando exigido, e associe a conta a uma profissão sorteada. \\
\hline
RF03 & Consulta de saldo e extrato & Apresentar o saldo da conta e o histórico de movimentações simuladas. \\
\hline
RF04 & Transferência entre contas & Permitir que o participante identifique uma conta de destino e transfira para ela um valor fictício. \\
\hline
RF05 & Pagamento de cobranças & Permitir a consulta de cobranças pendentes e seu pagamento com saldo em conta ou cartão, conforme as condições da operação. \\
\hline
RF06 & Cartões e faturas & Permitir a solicitação de cartão e a consulta de limite e fatura, além do pagamento da fatura no ambiente simulado. \\
\hline
RF07 & Investimentos & Permitir a aplicação e a retirada de valores fictícios, com consulta aos rendimentos apresentados pela aplicação. \\
\hline
RF08 & Compras no shopping & Permitir a consulta de produtos e a realização de compras fictícias com saldo ou cartão. \\
\hline
RF09 & Simulação de empréstimo & Permitir a estimativa de prestações a partir do valor solicitado e do prazo. O módulo possui interface e rotina de cálculo no navegador, cujo funcionamento precisa ser verificado em execução. \\
\hline
RF10 & Gerência de agência & Permitir a configuração da agência, a administração das contas vinculadas e o processamento de ciclos, incluindo o crédito de salário fictício. \\
\hline
\end{xltabular}

\section{Requisitos não funcionais}
\label{cap-concepcao-requisitos-nao-funcionais}

\section{Regras de negócio}
\label{cap-concepcao-regras}

\subsection{Contas, saldo e extrato}
\label{cap-concepcao-contas}

\subsection{Pagamentos e transferências simuladas}
\label{cap-concepcao-pagamentos}

\subsection{Simulação de crédito}
\label{cap-concepcao-credito}

\chapter{Desenvolvimento da Aplicação}
\label{cap-aplicacao}
\section{Arquitetura da aplicação}
\label{cap-desenvolvimento-arquitetura}

\section{Tecnologias utilizadas e justificativas}
\label{cap-desenvolvimento-tecnologias}

\section{Modelagem do banco de dados}
\label{cap-desenvolvimento-banco-dados}

\section{Implementação do back-end}
\label{cap-desenvolvimento-backend}

\subsection{Rotas da API e validação das operações}
\label{cap-desenvolvimento-rotas-api}

\subsection{Registro das movimentações e atualização do saldo}
\label{cap-desenvolvimento-movimentacoes-saldo}

\subsection{Implementação dos cálculos financeiros}
\label{cap-desenvolvimento-calculos-financeiros}

\section{Implementação do front-end}
\label{cap-desenvolvimento-frontend}

\subsection{Painel da conta e consulta ao extrato}
\label{cap-desenvolvimento-painel-conta}

\subsection{Telas de pagamento e transferência}
\label{cap-desenvolvimento-telas-pagamento}

\subsection{Tela de simulação de crédito}
\label{cap-desenvolvimento-tela-credito}

\section{Repositório e instruções de execução}
\label{cap-desenvolvimento-repositorio}


\chapter{Verificação e Resultados}
\label{cap-resultados}
\section{Cenários e procedimentos de teste}
\label{cap-verificacao-cenarios-testes}

\section{Verificação das movimentações, dos saldos e do extrato}
\label{cap-verificacao-movimentacoes-saldos-extrato}

\section{Verificação dos cálculos financeiros}
\label{cap-verificacao-calculos-financeiros}

\section{Resultados dos testes}
\label{cap-verificacao-resultados-testes}


\chapter{Considerações Finais}
\label{cap-consideracoes}
\section{Síntese do que foi desenvolvido}
\label{cap-consideracoes-sintese}

\section{Limitações e trabalhos futuros}
\label{cap-consideracoes-limitacoes-futuros}



% Introducao
% \input{Textual/cap1}

% Fundamentação Teórica
% \input{Textual/cap2}

% Revisão Bibliográfica
% \input{Textual/cap3}

% Metodologia
% \input{Textual/cap4}

% Considerações Finais
% \input{Textual/capConclusao}




% ---
% Finaliza a parte no bookmark do PDF, para que se inicie o bookmark na raiz
% ---
\bookmarksetup{startatroot}% 
% ---

% ---------------------------------------------------------------------------------------------
% ELEMENTOS PÓS-TEXTUAIS
% ---------------------------------------------------------------------------------------------
\postextual


% ---------------------------------------------------------------------------------------------
% Referências bibliográficas




% ---------------------------------------------------------------------------------------------
\bibliography{referenciais.bib}

% ---------------------------------------------------------------------------------------------
% Glossário
% ---------------------------------------------------------------------------------------------
%
% Consulte o manual da classe abntex2 para orientações sobre o glossário.
%
%\glossary

% ---------------------------------------------------------------------------------------------
% Apêndices
% ---------------------------------------------------------------------------------------------

% ---
% Inicia os apêndices
% ---
% \begin{apendicesenv}
%\partapendices
%\chapter{Modelagem do Sistema DOR}











%\end{apendicesenv}
% ---


% ---------------------------------------------------------------------------------------------
% Anexos
% ---------------------------------------------------------------------------------------------

% ---
% Inicia os anexos
% ---
%\begin{anexosenv}
%\partanexos
%\chapter{Protocolo de Aquisição de Imagens}








%\end{anexosenv}

% ---------------------------------------------------------------------------------------------
% INDICE REMISSIVO
% ---------------------------------------------------------------------------------------------

\printindex

\end{document}
